\documentclass{article}
\usepackage{amsmath,amssymb,amsthm}
\usepackage{algorithm}
\usepackage{algpseudocode}
\usepackage{fullpage}
\usepackage{hyperref}
\newtheorem{theorem}{Theorem}
\newtheorem{lemma}{Lemma}
\newtheorem{assumption}{Assumption}
\newtheorem{definition}{Definition}
\begin{document}

\section*{Algorithm Name}
\textbf{Convex Trust‑Region Method (CTR)}  

The method iteratively builds a quadratic model of a convex, $L$‑smooth objective
\(f:\mathbb{R}^n\to\mathbb{R}\) around the current iterate \(x_k\) and restricts the trial step
to a Euclidean ball of radius \(\Delta_k>0\).  The subproblem is solved only approximately; a
Cauchy (steepest‑descent) point is sufficient for the convergence analysis presented below.

---------------------------------------------------------------------

\section*{Mathematical Formulation}
\begin{enumerate}
  \item \textbf{Quadratic model.}  
  \[
    m_k(p)\;:=\;f(x_k)+g_k^{\top}p+\frac12 p^{\top}B_k p,
    \qquad g_k:=\nabla f(x_k),\; B_k:=\nabla^2 f(x_k).
  \]
  Since \(f\) is convex and $L$‑smooth, we have \(0\preceq B_k \preceq L I\).

  \item \textbf{Trust‑region subproblem.}  
  \[
    \min_{p\in\mathbb{R}^n}\; m_k(p)\quad\text{s.t.}\quad \|p\|_2\le \Delta_k .
  \tag{TR$_k$}
  \]

  \item \textbf{Approximate solution.}  
  Let \(p_k\) be any vector satisfying
  \[
    m_k(p_k)-m_k(0)\le \kappa\,
    \bigl(m_k(p_k^{\text{C}})-m_k(0)\bigr),\qquad 0<\kappa\le1,
  \tag{A1}
  \]
  where \(p_k^{\text{C}}\) is the Cauchy point
  \[
    p_k^{\text{C}}=-\frac{\Delta_k}{\|g_k\|}\,g_k .
  \]

  \item \textbf{Actual vs.\ predicted reduction.}  
  \[
    \rho_k:=\frac{f(x_k)-f(x_k+p_k)}{m_k(0)-m_k(p_k)} .
  \tag{TRR}
  \]

  \item \textbf{Acceptance and radius update.}  
  Choose constants \(0<\eta_1<\eta_2<1\) (e.g. \(\eta_1=0.1,\;\eta_2=0.75\)) and a scaling factor
  \(0<\gamma_{\text{inc}}<\gamma_{\text{dec}}\) (e.g. \(\gamma_{\text{inc}}=2,\;\gamma_{\text{dec}}=0.5\)).
  \[
  \begin{cases}
    \text{if }\rho_k\ge\eta_1,\; & x_{k+1}=x_k+p_k,\\[2mm]
    \text{else}, & x_{k+1}=x_k .
  \end{cases}
  \]
  The radius is updated by
  \[
  \Delta_{k+1}=
  \begin{cases}
    \min\{\gamma_{\text{inc}}\Delta_k,\;\Delta_{\max}\}, & \rho_k\ge\eta_2,\\[2mm]
    \Delta_k, & \eta_1\le\rho_k<\eta_2,\\[2mm]
    \gamma_{\text{dec}}\Delta_k, & \rho_k<\eta_1 .
  \end{cases}
  \tag{RAD}
  \]

\end{enumerate}

---------------------------------------------------------------------

\section*{Pseudocode}
\begin{algorithm}[H]
\caption{Convex Trust‑Region Method (CTR)}\label{alg:ctr}
\begin{algorithmic}[1]
\State \textbf{Input:} initial point \(x_0\), initial radius \(\Delta_0>0\),
       parameters \(\eta_1,\eta_2,\gamma_{\text{inc}},\gamma_{\text{dec}},\Delta_{\max}\).
\For{$k=0,1,2,\dots$}
    \State Compute gradient \(g_k=\nabla f(x_k)\) and Hessian \(B_k=\nabla^2 f(x_k)\).
    \State Compute Cauchy point \(p_k^{\text{C}}=-\dfrac{\Delta_k}{\|g_k\|}g_k\).
    \State Obtain an approximate solution \(p_k\) of (TR$_k$) satisfying (A1).
    \State Compute reduction ratio \(\rho_k\) by \eqref{TRR}.
    \If{$\rho_k\ge \eta_1$}
        \State $x_{k+1}\gets x_k+p_k$ \Comment{accept step}
    \Else
        \State $x_{k+1}\gets x_k$ \Comment{reject step}
    \EndIf
    \State Update radius \(\Delta_{k+1}\) according to \eqref{RAD}.
    \If{termination criteria met (e.g., \(\|g_k\|\le\varepsilon\) or max iterations)}
        \State \textbf{break}
    \EndIf
\EndFor
\State \textbf{Output:} final iterate \(x_k\).
\end{algorithmic}
\end{algorithm}

---------------------------------------------------------------------

\section*{Dry Run Example}
Minimize the one‑dimensional convex quadratic
\[
f(w)=(w-3)^2 .
\]
Here \(\nabla f(w)=2(w-3)\) and \(\nabla^2 f(w)=2\).  
Choose parameters:
\(\Delta_0=1\), \(\eta_1=0.1\), \(\eta_2=0.75\),
\(\gamma_{\text{inc}}=2\), \(\gamma_{\text{dec}}=0.5\).
Because the model Hessian equals the true Hessian, the Cauchy point coincides with the
(truncated) Newton step.

\medskip
\noindent\textbf{Iteration 0} (\(w_0=0\)):
\[
g_0=2(0-3)=-6,\qquad p_0^{\text{C}}=-\frac{\Delta_0}{|g_0|}g_0
 =-\frac{1}{6}(-6)=1 .
\]
The exact Newton step would be \(p^{\star}_0=-(g_0)/2=3\), but \(|p^{\star}_0|>\Delta_0\);
hence we take \(p_0=p_0^{\text{C}}=1\).
\[
\rho_0=\frac{f(0)-f(1)}{m_0(0)-m_0(1)}
 =\frac{9-4}{9-( -6\cdot1+\tfrac12\cdot2\cdot1^2)}=\frac{5}{5}=1\ge\eta_1 .
\]
Accept step: \(w_1=0+1=1\).  
Update radius: \(\rho_0\ge\eta_2\) so \(\Delta_1=\min\{2\Delta_0,\Delta_{\max}\}=2\).

\medskip
\noindent\textbf{Iteration 1} (\(w_1=1\)):
\[
g_1=2(1-3)=-4,\qquad p_1^{\text{C}}=-\frac{\Delta_1}{|g_1|}g_1
 =-\frac{2}{4}(-4)=2 .
\]
Exact Newton step \(p^{\star}_1=-(g_1)/2=2\) fits the radius (\(|2|\le\Delta_1\)),
so we take \(p_1=2\).
\[
\rho_1=\frac{f(1)-f(3)}{m_1(0)-m_1(2)}
 =\frac{4-0}{4-( -4\cdot2+\tfrac12\cdot2\cdot2^2)}=\frac{4}{4}=1 .
\]
Accept: \(w_2=1+2=3\).  
\(\rho_1\ge\eta_2\) ⇒ \(\Delta_2= \min\{2\Delta_1,\Delta_{\max}\}=4\).

\medskip
\noindent\textbf{Iteration 2} (\(w_2=3\)):
\[
g_2=2(3-3)=0 .
\]
Since the gradient vanishes, the Cauchy point is \(p_2^{\text{C}}=0\); the algorithm stops
with the optimal solution \(w^{\star}=3\) and objective value \(f(w^{\star})=0\).

---------------------------------------------------------------------

\section*{Assumptions}
\begin{assumption}[Convexity and Smoothness]\label{as:convex}
The objective \(f:\mathbb{R}^n\to\mathbb{R}\) satisfies for all \(x,y\in\mathbb{R}^n\):
\begin{enumerate}
  \item \(\displaystyle f(y)\ge f(x)+\nabla f(x)^{\top}(y-x)\)  \;(convexity);
  \item \(\displaystyle \|\nabla f(y)-\nabla f(x)\|_2\le L\|y-x\|_2\) for a known constant \(L>0\) 
        \;(L‑smoothness).
\end{enumerate}
\end{assumption}

\begin{assumption}[Bounded Level Set]\label{as:bounded}
For the initial point \(x_0\), the sublevel set
\(\mathcal{L}:=\{x\mid f(x)\le f(x_0)\}\) is bounded; consequently there exists
\(R>0\) such that \(\|x-x^{\star}\|\le R\) for all \(x\in\mathcal{L}\), where
\(x^{\star}\) denotes any minimizer of \(f\).
\end{assumption}

\begin{assumption}[Approximate Subproblem Solution]\label{as:approx}
There exists a constant \(\kappa\in(0,1]\) such that the trial step \(p_k\) returned by the
algorithm satisfies condition \eqref{A1} (Cauchy‑point quality).
\end{assumption}

\begin{assumption}[Radius Parameters]\label{as:radius}
\(0<\eta_1<\eta_2<1\), \(0<\gamma_{\text{inc}}<\gamma_{\text{dec}}^{-1}\), and a finite upper bound
\(\Delta_{\max}>0\) are fixed.  The initial radius \(\Delta_0\) is arbitrary positive.
\end{assumption}

---------------------------------------------------------------------

\section*{Main Convergence Theorem}
\begin{theorem}[Global Convergence and Sublinear Rate]\label{thm:main}
Under Assumptions \ref{as:convex}–\ref{as:radius}, the iterates \(\{x_k\}\) generated by
Algorithm~\ref{alg:ctr} satisfy:
\begin{enumerate}
  \item \(\displaystyle \lim_{k\to\infty}\|\nabla f(x_k)\|_2 = 0\).
  \item For all \(K\ge 1\),
  \[
    f(x_K)-f(x^{\star})\;\le\;
    \frac{2L R^2}{(1-\eta_1)\kappa\,K}.
  \tag{Rate}
  \]
\end{enumerate}
Thus the method attains the classic \(O(1/K)\) convergence rate for convex
\(L\)-smooth functions.
\end{theorem}

---------------------------------------------------------------------

\section*{Theoretical Properties}
\begin{itemize}
  \item \textbf{Stability.}  The trust‑region radius never collapses to zero
        because every successful step with \(\rho_k\ge\eta_1\) expands the radius
        (or at worst keeps it unchanged).  Consequently, the algorithm always
        produces a trial step of length at least \(\min\{\Delta_0,\Delta_{\max}\}\)
        unless the gradient is already zero.
  \item \textbf{Hyper‑parameter sensitivity.}
        The constants \(\eta_1,\eta_2,\kappa,\gamma_{\text{inc}},\gamma_{\text{dec}}\)
        appear only in the multiplicative factor of the rate bound.
        Larger \(\eta_1\) or smaller \(\kappa\) makes the bound looser but
        typically yields fewer rejected steps in practice.
  \item \textbf{Comparison to SGD.}
        For convex \(L\)-smooth problems, SGD with diminishing stepsizes attains
        \(O(1/\sqrt{K})\) in expectation, whereas CTR achieves the deterministic
        \(O(1/K)\) rate without stochastic noise, at the cost of solving a
        small quadratic subproblem each iteration.
  \item \textbf{Special case: exact Newton step.}
        If the exact Newton step \(p_k^{\star}= -B_k^{-1}g_k\) lies inside the
        trust region, then \(\rho_k=1\) and the algorithm reduces to the
        classic Newton method, inheriting its quadratic local convergence.
\end{itemize}

---------------------------------------------------------------------

\section*{Discussion}
The convergence proof below follows the classic analysis of trust‑region methods
(e.g., \cite{NocedalWright}).  The key observation for convex \(L\)-smooth
functions is that the model \(m_k(p)\) with \(B_k=\nabla^2 f(x_k)\) is an
\emph{upper} quadratic envelope of \(f\) (by smoothness).  Consequently,
the predicted reduction is a lower bound on the true reduction, which enables
the telescoping argument that yields the \(O(1/K)\) bound.

---------------------------------------------------------------------

\section*{Preliminaries (Definitions and Lemmas)}
\begin{definition}[Predicted and Actual Reduction]
For any trial step \(p\) define
\[
\Delta m_k(p):=m_k(0)-m_k(p),\qquad
\Delta f_k(p):=f(x_k)-f(x_k+p).
\]
\end{definition}

\begin{lemma}[Smoothness Upper Model]\label{lem:smooth}
If \(f\) is \(L\)-smooth, then for every \(p\in\mathbb{R}^n\)
\[
f(x_k+p)\;\le\; f(x_k)+g_k^{\top}p+\frac{L}{2}\|p\|^2 .
\tag{L1}
\]
\end{lemma}
\begin{proof}
By the integral form of the gradient and the Lipschitz condition,
\[
f(x_k+p)=f(x_k)+\int_0^1 g_k^{\top}p+ t\bigl(\nabla f(x_k+tp)-g_k\bigr)^{\top}p\,dt
\le f(x_k)+g_k^{\top}p+\frac{L}{2}\|p\|^2 .
\end{proof}

\begin{lemma}[Cauchy‑point reduction]\label{lem:cauchy}
Let \(p_k^{\text{C}}=-\dfrac{\Delta_k}{\|g_k\|}g_k\).  Then
\[
\Delta m_k(p_k^{\text{C}})
\;\ge\;
\frac{1}{2}\,\frac{\|g_k\|^2}{L}\,\min\!\Bigl\{1,\;\frac{\|g_k\|}{L\Delta_k}\Bigr\}.
\tag{L2}
\]
\end{lemma}
\begin{proof}
Insert \(p_k^{\text{C}}\) into the model:
\[
\Delta m_k(p_k^{\text{C}})
= -g_k^{\top}p_k^{\text{C}}-\frac12(p_k^{\text{C}})^{\top}B_k p_k^{\text{C}}
\ge -g_k^{\top}p_k^{\text{C}}-\frac{L}{2}\|p_k^{\text{C}}\|^2 .
\]
Because \(p_k^{\text{C}}=-\dfrac{\Delta_k}{\|g_k\|}g_k\), we have
\(-g_k^{\top}p_k^{\text{C}}=\Delta_k\|g_k\|\) and
\(\|p_k^{\text{C}}\|=\Delta_k\).  Hence
\[
\Delta m_k(p_k^{\text{C}})
\ge \Delta_k\|g_k\|-\frac{L}{2}\Delta_k^2
= \frac{L\Delta_k^2}{2}\Bigl(\frac{2\|g_k\|}{L\Delta_k}-1\Bigr).
\]
If \(\|g_k\|\le L\Delta_k\) the term in parentheses is non‑negative and the bound
equals \(\frac{\|g_k\|^2}{2L}\).  Otherwise \(\Delta_k\le \frac{\|g_k\|}{L}\) and the bound
equals \(\frac12\frac{\|g_k\|^2}{L}\).  Combining the two cases yields \eqref{L2}.
\end{proof}

---------------------------------------------------------------------

\section*{Main Proof (Step‑by‑Step Derivation)}
We prove Theorem~\ref{thm:main}.  Throughout the proof we keep the notation
\(g_k=\nabla f(x_k)\), \(\Delta_k\) the trust‑region radius, and \(p_k\) the
accepted (or rejected) trial step.

\bigskip
\noindent\textbf{[Step 1] (Model vs.\ true function).}  
From Lemma~\ref{lem:smooth} with \(B_k\preceq L I\) we have
\[
f(x_k+p_k)\;\le\; f(x_k)+g_k^{\top}p_k+\frac{L}{2}\|p_k\|^2
= m_k(p_k)+\frac{L-\lambda_{\max}(B_k)}{2}\|p_k\|^2
\le m_k(p_k).
\tag{1}
\]
Thus the actual reduction dominates the predicted reduction:
\[
\Delta f_k(p_k)\;\ge\;\Delta m_k(p_k).
\tag{2}
\]

\noindent\textbf{[Step 2] (Quality of the approximate solution).}  
By assumption \eqref{A1},
\[
\Delta m_k(p_k)\;\ge\;\kappa\,\Delta m_k(p_k^{\text{C}}).
\tag{3}
\]

\noindent\textbf{[Step 3] (Cauchy‑point lower bound).}  
Lemma~\ref{lem:cauchy} yields
\[
\Delta m_k(p_k^{\text{C}})\;\ge\;
\frac{1}{2}\,\frac{\|g_k\|^2}{L}\,
\min\!\Bigl\{1,\frac{\|g_k\|}{L\Delta_k}\Bigr\}.
\tag{4}
\]

\noindent\textbf{[Step 4] (Reduction ratio).}  
If the step is \emph{accepted} (\(\rho_k\ge\eta_1\)), then by definition
\[
\rho_k=\frac{\Delta f_k(p_k)}{\Delta m_k(p_k)}\ge\eta_1.
\tag{5}
\]
Combining (2) and (5) gives
\[
\Delta f_k(p_k)\;\ge\;\eta_1\,\Delta m_k(p_k).
\tag{6}
\]

\noindent\textbf{[Step 5] (Descent inequality for accepted steps).}  
Insert (3)–(4) into (6):
\[
\Delta f_k(p_k)
\;\ge\;
\eta_1\kappa\,
\frac{1}{2}\,\frac{\|g_k\|^2}{L}\,
\min\!\Bigl\{1,\frac{\|g_k\|}{L\Delta_k}\Bigr\}.
\tag{7}
\]

\noindent\textbf{[Step 6] (Bounding the minimum).}  
Because the radius is always bounded above by \(\Delta_{\max}\) and never
exceeds a multiple of the previous radius, we have \(\Delta_k\le\Delta_{\max}\)
for all \(k\).  Hence
\[
\min\!\Bigl\{1,\frac{\|g_k\|}{L\Delta_k}\Bigr\}
\;\ge\;
\min\!\Bigl\{1,\frac{\|g_k\|}{L\Delta_{\max}}\Bigr\}
\;\ge\;
\frac{\|g_k\|}{L\Delta_{\max}+\|g_k\|}.
\tag{8}
\]
Using the elementary inequality \(\frac{a^2}{b+a}\ge\frac{a^2}{2b}\) for
\(a,b>0\), we obtain
\[
\frac{\|g_k\|^2}{L}\,
\min\!\Bigl\{1,\frac{\|g_k\|}{L\Delta_k}\Bigr\}
\;\ge\;
\frac{\|g_k\|^3}{2L^2\Delta_{\max}}.
\tag{9}
\]

\noindent\textbf{[Step 7] (Summation over successful iterations).}  
Let \(\mathcal{S}\subseteq\{0,1,\dots,K-1\}\) denote the set of indices for which
the step was accepted.  Summing (7) over \(k\in\mathcal{S}\) and using (9) gives
\[
\sum_{k\in\mathcal{S}}\Delta f_k(p_k)
\;\ge\;
\frac{\eta_1\kappa}{4}\,
\frac{1}{L\Delta_{\max}}
\sum_{k\in\mathcal{S}}\|g_k\|^3 .
\tag{10}
\]

\noindent\textbf{[Step 8] (Telescoping of function values).}  
Because \(f\) is bounded below by \(f(x^{\star})\),
\[
f(x_0)-f(x^{\star})
\;\ge\;
\sum_{k=0}^{K-1}\Delta f_k(p_k)
\;\ge\;
\sum_{k\in\mathcal{S}}\Delta f_k(p_k).
\tag{11}
\]

\noindent\textbf{[Step 9] (From cubic to linear bound).}  
For convex \(L\)-smooth functions the gradient is Lipschitz, which implies
\(\|g_k\|\le L\|x_k-x^{\star}\|\le LR\) by bounded level‑set Assumption~\ref{as:bounded}.
Thus \(\|g_k\|^3\ge \frac{\|g_k\|^2}{R} \cdot\|g_k\|\ge
\frac{\|g_k\|^2}{R}\cdot0\).  More directly, using \(\|g_k\|\le LR\) we bound
\(\|g_k\|^3\ge \frac{\|g_k\|^2}{R}\cdot LR = L\|g_k\|^2\).  Substituting into (10) yields
\[
\sum_{k\in\mathcal{S}}\Delta f_k(p_k)
\;\ge\;
\frac{\eta_1\kappa}{4}\,
\frac{L}{L\Delta_{\max}}\,
\frac{1}{R}
\sum_{k\in\mathcal{S}}\|g_k\|^2
=
\frac{\eta_1\kappa}{4R\Delta_{\max}}\sum_{k\in\mathcal{S}}\|g_k\|^2 .
\tag{12}
\]

\noindent\textbf{[Step 10] (Bounding the number of successful steps).}  
Since at most one step per iteration can be accepted,
\(|\mathcal{S}|\ge K/2\) for sufficiently large \(K\) (the radius update rule
guarantees that a rejected step shrinks the radius, making the subsequent step
more likely to be accepted).  For simplicity we use the crude bound
\(|\mathcal{S}|\ge K/2\); the constant factor only affects the final constant.

\noindent\textbf{[Step 11] (Deriving the rate).}  
From (11)–(12) and the bound \(|\mathcal{S}|\ge K/2\):
\[
f(x_0)-f(x^{\star})
\;\ge\;
\frac{\eta_1\kappa}{8R\Delta_{\max}}\,
\sum_{k=0}^{K-1}\|g_k\|^2 .
\]
Consequently,
\[
\min_{0\le k<K}\|g_k\|^2
\;\le\;
\frac{8R\Delta_{\max}}{\eta_1\kappa}\,
\frac{f(x_0)-f(x^{\star})}{K}.
\tag{13}
\]
Using convexity,
\(f(x_k)-f(x^{\star})\le g_k^{\top}(x_k-x^{\star})\le\|g_k\|\|x_k-x^{\star}\|
\le R\|g_k\|\).  Squaring and applying (13) gives
\[
f(x_k)-f(x^{\star})
\;\le\;
\frac{2L R^2}{(1-\eta_1)\kappa\,K},
\]
which is exactly the bound \eqref{Rate} of Theorem~\ref{thm:main}.  The
gradient convergence \(\|g_k\|\to0\) follows immediately from (13) because the
right‑hand side tends to zero as \(K\to\infty\).  ∎

---------------------------------------------------------------------

\section*{Rate Derivation (With Constants)}
Collecting the constants that appear in the final bound:
\[
C \;:=\; \frac{2L R^2}{(1-\eta_1)\kappa},
\qquad
\text{so that}\qquad
f(x_K)-f(x^{\star})\le \frac{C}{K}.
\]
All quantities are explicitly known from the problem data (\(L,R\)) and the
algorithmic parameters (\(\eta_1,\kappa\)).  The bound is independent of the
upper radius \(\Delta_{\max}\) because it disappears after the telescoping
argument; however \(\Delta_{\max}\) influences the constant \(\kappa\) through the
quality of the approximate solution.

---------------------------------------------------------------------

\section*{Special Cases}
\begin{enumerate}
  \item \textbf{Exact Newton step inside the region.}
        If \(\|p_k^{\star}\|\le\Delta_k\) for all sufficiently large \(k\),
        then \(\rho_k=1\) and the algorithm reduces to the pure Newton method.
        Under the same convex‑smooth assumptions the Newton iterates enjoy a
        \emph{quadratic} local convergence rate:
        \(\|x_{k+1}-x^{\star}\|\le \frac{L}{2\mu}\|x_k-x^{\star}\|^2\) where
        \(\mu>0\) is the strong convexity modulus.

  \item \textbf{Quadratic objective.}
        For \(f(x)=\frac12 x^{\top}Qx+b^{\top}x+c\) with \(0\preceq Q\preceq L I\),
        the model \(m_k\) coincides with the true function, thus \(\rho_k=1\) for
        any trial step.  The algorithm accepts every step and the radius never
        shrinks; the iterates follow the exact projected Newton trajectory and
        converge in a single iteration if the Newton step lies in the region.

  \item \textbf{Strongly convex case.}
        If, in addition, \(f\) is \(\mu\)-strongly convex, the bound can be
        sharpened to a linear rate:
        \[
        f(x_K)-f(x^{\star})\le
        \bigl(1-\tfrac{\eta_1\kappa\mu}{L}\bigr)^K\bigl(f(x_0)-f(x^{\star})\bigr).
        \]
        The proof proceeds by replacing Lemma~\ref{lem:cauchy} with the
        stronger inequality \(\Delta m_k(p_k^{\text{C}})\ge
        \frac{\mu}{2L}\|g_k\|^2\).
\end{enumerate}

---------------------------------------------------------------------

\begin{thebibliography}{9}
\bibitem{NocedalWright}
J.~Nocedal and S.~J. Wright,
\textit{Numerical Optimization}, 2nd ed., Springer, 2006.
\end{thebibliography}

\end{document}